\chapter{CFFM-Net}\label{ch:cffm}

\section{The whole network}

\method{} detects objects from a visible and a thermal camera at once (Figure~\ref{fig:overall}). Each camera has
its own YOLO26-n backbone with separate weights; both start from the same COCO-pretrained weights, the thermal
image repeated to three channels. The two backbones run side by side, and at four pyramid levels, P2 to P5, their
feature maps are fused. At P2, stride 4, the fusion is a reliability-weighted average, because a scan over a
$160\times160$ map buys little for its cost. At P3, P4 and P5 it is the cross-modal fusion Mamba block, CMFM. The
fused maps $Z_2$ to $Z_5$ then enter an unchanged YOLO26 neck extended with a stride-4 path, and four heads predict
boxes at strides 4, 8, 16 and 32.

\widefig{fig_cffm_overall}{\method{} as trained in the pilot, drawn with the real data of one LLVIP test pair. The
feature maps are the channel-averaged activations of each backbone stage; the small maps $r^{\theta}$ are the
thermal reliability the model assigned at each fusion level. 6.03 million parameters: 1.37 million in each
backbone, 2.15 million in the fusion blocks and the rest in the neck and heads.}{fig:overall}

The fusion never feeds back into the backbones. Each stream keeps its own features all the way to P5, so one
camera's failure cannot corrupt the other's representation, and the fusion blocks decide how much of each to pass
to the neck. Figure~\ref{fig:compact} gives the same network in compact form.

\fig[0.82\textwidth]{fig_architecture}{\method{} in compact form: one backbone stage per stream, the fusion at each
level, the shared neck and the heads. The block repeats four times, once per pyramid level.}{fig:compact}

\section{Inside the fusion block}

Figure~\ref{fig:cmfm} opens CMFM. A small reliability head looks at both feature maps and at their difference and
outputs two maps, $r^{v}$ and $r^{\theta}$, between zero and one: how far the block should trust each camera at
each location. The thermal features are first shifted by a learned offset field of at most four cells, which
corrects small misregistration between the cameras. Both streams are projected to half their channels, then
interleaved token by token, $x^{v}_{1}, x^{\theta}_{1}, x^{v}_{2}, x^{\theta}_{2}, \dots$, so that the scan sees both
cameras at the same place one after the other. The gated scan runs over this sequence in four directions, its
output is merged and normalised, joined with a local depthwise branch, and projected back. Because that last
projection starts at zero, a new block computes exactly the reliability-weighted average of its inputs, and
training can only add to it.

\fig[0.66\textwidth]{fig_cmfm}{The cross-modal fusion Mamba block. The reliability maps are used twice: they gate
the scan and they weight the residual average. The output projection $W_o$ starts at zero.}{fig:cmfm}

\section{The gated recurrence}

The single change that makes the scan selective between cameras is drawn in Figure~\ref{fig:gated}. Each token's
step size $\Delta_k$ is multiplied by the reliability $r_k$ of the camera that produced it. When $r_k$ approaches
zero, $\Delta_k$ does too, so the decay approaches one and the input term zero: the token neither writes into the
state nor erases it. A dazzled or missing camera is skipped, and the memory built from the other camera carries
on. The unit tests check this property directly: with thermal reliability set to zero, changing the thermal
features does not change the visible outputs.

\fig[0.66\textwidth]{fig_gated_scan}{One step of the reliability-gated selective scan. The reliability $r_k$ scales
the step size, which is the knob Mamba uses to decide what to keep.}{fig:gated}

\section{The numbers}

\begin{tabularx}{\textwidth}{@{}P{5.2cm}LLL@{}}
\hd{model} & \hd{parameters} & \hd{GFLOPs} & \hd{latency, T4 FP16} \\
YOLO26-n, one camera & 2.50 M & 5.9 & 17.0 ms \\
two-stream concat & 4.07 M & 9.7 & 23.4 ms \\
\method{}, pilot & 6.03 M & 18.4 & 41.9 ms \\
\end{tabularx}

Latency is measured at batch 1 on one T4 with the PyTorch model, before any engine optimisation. \method{} costs
1.79 times the concat baseline, inside the limit of two set before training.

\section{Versions}

The pilot version uses YOLO26's stride-4 head. A head-control experiment showed that this head, not the fusion,
cost about 1.8 AP on LLVIP, so the second version, \method{} v2, returns to the standard three-level head. It also
trains with modality dropout, blacking out one camera at random so the model learns to work without it, and keeps
the fusion weights in single precision with a floor on their sum, which removed a numerical overflow. In the first
clean comparison on M3FD, \method{} v2 reached 0.488 AP against 0.483 for the concat baseline trained the same way.
